\documentclass[b5paper,10pt,onecolumn,twoside,fontset=windows]{ctexart}
\pagestyle{empty}%不显示页码
\usepackage[total={130mm,203mm}, top=20.3mm, headheight=20pt, headsep=4truemm]{geometry}%调整版式
\usepackage{tabularx}
\usepackage{amssymb}
\setmainfont{Times New Roman}

% \renewcommand{\large}{\fontsize{11}{14}\selectfont}%调整字体大小
\newcommand{\vdistance}{\\[4mm]}%调整行间隔

\begin{document}
\begin{flushright}
\zihao{-4}
\textbf{Studies in Logic\phantom{|}\\
Volume 19, Number 5, October 2026\phantom{|}}
\end{flushright}
\vspace{0.5ex} {\rmfamily{\bfseries TABLE OF CONTENTS}} \vspace{-0.5ex}

\noindent\hrulefill{\rmfamily} \vspace{2ex}\\
\begin{tabularx}{\textwidth}{l@{\hspace{4em}}r}

%中间主体

\large On the Example Paradox: A Study Based on the Extended X Criterion\\
\textit{}
& 1\\

\end{tabularx}
\end{document}
